% ==================== CHAPTER 2 ====================
\begin{Arabic}
\chapter{الإطار النظري والدراسات السابقة}
\clearpage

\section{أنظمة معالجة الإيماءات في تطبيقات الهواتف المحمولة}

شهد مجال تطوير تطبيقات الهواتف الذكية تطوراً متسارعاً في التعامل مع إيماءات اللمس، إذ انتقل من نماذج بدائية تعتمد على الأحداث المنفردة إلى أطر عمل متكاملة قادرة على التعرف على أنماط حركية معقدة في الوقت الفعلي. تُعدّ إيماءات اللمس (Touch Gestures) اليوم العمود الفقري لتجربة المستخدم في أي تطبيق محمول، إذ تُشير الدراسات إلى أن ما يزيد على 70\% من تفاعلات المستخدم مع تطبيقات الهواتف تتم عبر إيماءات لمسية متنوعة.

يُقدم \textenglish{Dabit} \cite{dabit2019} في كتابه \textenglish{React Native in Action} أحد أشمل المراجع التقنية في هذا المجال، إذ يُفصّل آلية عمل نظام الاستجابة للمس (\textenglish{Gesture Responder System}) في \textenglish{React Native}. ويُبيّن أن الإطار الافتراضي المُقدَّم (\textenglish{PanResponder}) يوفر واجهة للتعامل مع أحداث اللمس الأساسية، غير أن هذا النظام يعاني من قيد بنيوي جوهري يتمثل في تشغيله الكامل داخل مسار \textenglish{JavaScript}، مما يُعرّضه للتأثر بأي ضغط حسابي على هذا المسار. وقد أفرز هذا القيد حاجة ملحّة لأطر عمل أكثر كفاءة تُجري المعالجة على المسار الأصلي للواجهة مباشرةً.

\subsection{تطور أطر معالجة الإيماءات}

تتسم الأجيال الأولى من أطر الإيماءات في \textenglish{React Native} بمحدودية واضحة؛ فقد كانت تعتمد بشكل حصري على الجسر (\textenglish{The Bridge}) لنقل أحداث اللمس من الطبقة الأصلية إلى بيئة \textenglish{JavaScript}. وكما يُوضح \textenglish{Dabit} \cite{dabit2019}، فإن هذا الجسر يُسبب تأخيراً كامناً (\textenglish{Latency}) في استجابة الإيماءات بسبب التسلسل (\textenglish{Serialization}) والإرسال (\textenglish{Dispatch}) غير المتزامن للأحداث عبر حدود الخيوط (\textenglish{Thread Boundaries}).

في المقابل، تُرسي مكتبة \textenglish{react-native-gesture-handler} نهجاً مغايراً جذرياً، إذ تُنفّذ منطق التعرف على الإيماءات بالكامل داخل المسار الأصلي للواجهة (\textenglish{UI Thread})، مما يُلغي تأخير الجسر ويضمن استجابة آنية تعكس حركة الإصبع بدقة متناهية. هذا التحول المعماري يُمثّل نقلة نوعية في تصميم مكتبات الإيماءات للمنصات المتقاطعة.

\section{تصنيف الإيماءات اللمسية وخصائصها التقنية}

تنقسم الإيماءات اللمسية في سياق التطبيقات المحمولة إلى فئتين رئيسيتين تتباينان في طبيعة الأحداث التي تُولّدانها:

\textbf{أولاً: الإيماءات المنفصلة (\textenglish{Discrete Gestures})} وهي إيماءات تُنتج حدثاً واحداً محدداً اكتمال شروطه، وتشمل الضغطة الواحدة (\textenglish{Tap})، والضغطة المزدوجة (\textenglish{Double Tap})، والضغط المطوّل (\textenglish{Long Press}). تعتمد آلية التعرف عليها أساساً على عاملَي الزمن والمسافة؛ إذ تُعرَّف الضغطة الواحدة بأنها لمس لا تتجاوز مدته 500 ميلي ثانية مع تحقيق قيد المسافة (أقل من 20 بكسل).

\textbf{ثانياً: الإيماءات المستمرة (\textenglish{Continuous Gestures})} وهي إيماءات تُنتج تدفقاً متواصلاً من الأحداث طوال فترة تنفيذها، وتشمل السحب (\textenglish{Drag})، والتحريك الشامل (\textenglish{Pan})، والقرص (\textenglish{Pinch})، والتدوير (\textenglish{Rotation}). ويُؤكد \textenglish{Masiello} و\textenglish{Friedmann} \cite{masiello2017} أن معالجة هذه الإيماءات تستلزم تحديثاً مستمراً لعناصر الواجهة بمعدل لا يقل عن 60 إطاراً في الثانية لتوفير تجربة سلسة وطبيعية.

\subsection{المعمارية الداخلية لأنظمة التعرف على الإيماءات}

تعتمد أنظمة التعرف على الإيماءات الحديثة نموذج آلة الحالات المتناهية (\textenglish{Finite State Machine}) للتمييز بين مراحل الإيماءة المختلفة. تمر كل إيماءة في \textenglish{gesture-kit} عبر ست حالات متسلسلة: \textenglish{UNDETERMINED} (الحالة الأولية)، و\textenglish{BEGAN} (بداية اللمس)، و\textenglish{ACTIVE} (تجاوز حد التعرف)، و\textenglish{END} (اكتمال الإيماءة)، و\textenglish{FAILED} (عدم استيفاء الشروط)، و\textenglish{CANCELLED} (إلغاء خارجي). يُحقق هذا النموذج تمييزاً دقيقاً بين الإيماءات المتشابهة ويمنع التضارب بينها.

\section{التحديات التقنية في تطوير مكتبات الإيماءات متعددة المنصات}

\subsection{مشكلة تأخر الجسر والحلول الحديثة}

تُمثّل مشكلة تأخر الجسر (\textenglish{Bridge Latency}) التحدي الأكثر إلحاحاً في تطوير إيماءات \textenglish{React Native}. يُبيّن \textenglish{Dabit} \cite{dabit2019} أن الجسر يعمل بصورة غير متزامنة، مما يعني أن كل حدث لمس يستلزم دورة كاملة من التسلسل وإرسال الرسائل وإلغاء التسلسل قبل أن يصل إلى مُعالج \textenglish{JavaScript}. في سيناريوهات الإيماءات المستمرة كالسحب والتكبير، يُنتج اللمس مئات الأحداث في الثانية الواحدة، مما يُضاعف الضغط على هذا الجسر ويُفضي إلى تذبذب ملحوظ في معدل الإطارات.

تُعالج مكتبة \textenglish{react-native-reanimated} هذه الإشكالية عبر مفهوم القيم المشتركة (\textenglish{Shared Values})، وهي متغيرات يمكن قراءتها وكتابتها من كلا المسارَين (الأصلي و\textenglish{JavaScript}) دون الحاجة إلى عبور الجسر. كما تُتيح وظيفة \textenglish{runOnUI} تنفيذ أوامر برمجية مباشرة على مسار الواجهة، مما يمنح المطور تحكماً دقيقاً في الأداء.

\subsection{إدارة تضارب الإيماءات}

يُعدّ تضارب الإيماءات (\textenglish{Gesture Conflict}) من أكثر التحديات تعقيداً في تطوير واجهات اللمس المتعددة. يطرح \textenglish{Masiello} و\textenglish{Friedmann} \cite{masiello2017} نموذج التسلسل الهرمي (\textenglish{Gesture Priority Hierarchy}) كحل لهذه المعضلة، حيث تُمنح الأولوية للإيماءة الأكثر تخصصاً أو الأقرب للعنصر المستهدف. تُقدم مكتبة \textenglish{react-native-gesture-handler} آليتَي \textenglish{simultaneousWithExternalGesture} و\textenglish{requireExternalGestureToFail} لضبط علاقات الأولوية بين الإيماءات المتزامنة.

\section{مكتبات الإيماءات والأنماط المعمارية للمكتبات مفتوحة المصدر}

\subsection{مبادئ تصميم المكتبات القابلة لإعادة الاستخدام}

يُرسي \textenglish{Casciaro} و\textenglish{Mammino} \cite{casciaro2020} في كتابهما \textenglish{Node.js Design Patterns} أسس تصميم المكتبات الاحترافية القابلة للتوزيع. ومن أبرز المبادئ التي تنطبق مباشرة على تصميم \textenglish{gesture-kit}: مبدأ المسؤولية الواحدة (\textenglish{Single Responsibility Principle}) الذي يقتضي أن تُتقن المكتبة مهمة واحدة محددة، وهي معالجة الإيماءات. كذلك مبدأ الإفصاح الصريح عن التبعيات (\textenglish{Explicit Dependencies}) عبر آلية \textenglish{peerDependencies} في ملف \textenglish{package.json}، لضمان التوافق مع إصدارات \textenglish{React Native} المختلفة في مشاريع المطورين.

\subsection{نمط الخطافات في منظومة React}

استحدث \textenglish{React} نمط الخطافات (\textenglish{Hooks}) بوصفه الطريقة المعيارية للتعامل مع المنطق القابل لإعادة الاستخدام. في سياق مكتبات الإيماءات، يُتيح هذا النمط تغليف منطق التعرف على الإيماءة بأكمله داخل دالة واحدة تُعيد كائن الإيماءة (\textenglish{GestureObject}) جاهزاً للتمرير إلى مكوّن \textenglish{GestureDetector}. يُشير \textenglish{Casciaro} و\textenglish{Mammino} \cite{casciaro2020} إلى أن هذا النمط يُحقق توازناً مثالياً بين سهولة الاستخدام وإمكانية التركيب والتوسعة.

\section{نظم الحساسات الداخلية في الأجهزة المحمولة}

تُضيف الحساسات المدمجة في الأجهزة الذكية بعداً تفاعلياً إضافياً لم تستوعبه كثير من مكتبات الإيماءات التقليدية. يُوضح \textenglish{Masiello} و\textenglish{Friedmann} \cite{masiello2017} أن مستشعرات التقارب (\textenglish{Proximity Sensors}) ومقاييس التسارع (\textenglish{Accelerometers}) تتدفق بياناتها عبر طبقة تجريد العتاد (\textenglish{Hardware Abstraction Layer}) وصولاً إلى تطبيق \textenglish{React Native}. 

وقد بحث \textenglish{Wobbrock} وزملاؤه \cite{wobbrock2009} في تصميم الإيماءات المعرفة من قبل المستخدم، مؤكدين أن دمج مدخلات إضافية يعزز من قابلية الاستخدام وتنوع التفاعلات. وفي هذا السياق، قدم \textenglish{Liu} وزملاؤه \cite{liu2019} نظام \textenglish{uWave} الذي يستفيد من مقياس التسارع للتعرف على إيماءات شخصية تعتمد على حركة اليد، مما يثبت إمكانية استخدام الحساسات الفيزيائية كأداة إدخال رديفة للشاشة اللمسية.

إن الدمج الفعّال بين أحداث اللمس وبيانات الحساسات في إطار عمل موحد يُفتح آفاقاً واسعة أمام تطبيقات الأمان الذكي، والألعاب التفاعلية، وتطبيقات إمكانية الوصول.

\section{خلاصة الإطار النظري}

يتضح من مراجعة الأدبيات أن تطوير إطار عمل شامل للإيماءات يستدعي معالجة ثلاثة تحديات رئيسية متشابكة: أولها إزالة تأثير الجسر ونقل المعالجة إلى المسار الأصلي، وثانيها تقديم واجهة برمجية (\textenglish{API}) نظيفة ومتسقة تستوعب كافة فئات الإيماءات بدءاً من الأساسية وانتهاءً بالتحويلات الهندسية المعقدة، وثالثها ضمان التوافق مع معايير النشر المهني على منصة \textenglish{npm} لاعتماد المكتبة في مشاريع حقيقية. وقد أرست المراجع الثلاثة الأساسية المُدرجة في هذا البحث إطاراً مفاهيمياً متكاملاً يُرشد التصميم المعماري للمشروع ويُجيب عن الأسئلة التصميمية الجوهرية.

\end{Arabic}
